% Time form: Present tense

\chapter{Introduction}\label{ch:introduction}

According to the World Economic Forum’s Global Risks Report 2026, cyber insecurity is listed among the top ten global risks and is projected to rise significantly from position 30 to position 5 in the long term, highlighting the growing adverse threat of \gls{AI} technologies \cite{WEF2025GlobalRisks}.

Complementing this global perspective, the CrowdStrike Global Threat Report highlights a significant rise in AI-enabled adversaries \cite{CrowdStrike2025Report}.

The absence of, or understaffing of, security experts, paired with the emerging complexity of systems and business models \ cite {Figueredo2024RCVaR}, highlights the increasing need for automated tools and frameworks capable of identifying and mitigating emerging cyber threats.

To address this, the \textit{Network Swiss Army Knife (NSAK)}, a modular, open-source framework, is designed to orchestrate security scenarios consisting of reusable attack drills. The framework aims to support automated security testing and adversary emulation in controlled network environments.

The development of the NSAK framework builds upon previous project work, in which the core architecture, simplified configuration management, and a basic scenario execution model were developed.\cite{nsakRepository2026}.

This thesis further extends the framework by enhancing configuration management for operating the NSAK device, establishing a reproducible test environment, and investigating how agentic AI can support automated red-team and blue-team scenarios.

\todo{Match the structure overview to the actual structure}
\todo{sharpen research question and contribution}

The thesis is structured as follows: \textbf{Chapter 2 }reviews related work and analyzes how selected security frameworks integrate \gls{AI} into their products. \textbf{Chapter 3} introduces the fundamental concepts and outlines the key design decisions that form the basis of this work. \textbf{Chapter 4 }describes the research methodology, including the approach used for evaluation and data collection. \textbf{Chapter 5 }presents the implementation of the proposed solution, followed by \textbf{Chapter 6}, which evaluates its effectiveness in a virtual and physical test environments. Finally, \textbf{Chapter 7} discusses the results and their implications, while \textbf{Chapter 8} concludes the thesis and highlights potential directions for future work.
